% !PRINTED none
% Abschnitt-title page, no running folio visible in scan: "Twenty-first Chapter. The Norms of the Class Invariants f(omega)." heading, rule, then § 140 heading and body text begin.

\begin{center}
\gesperrt{Twenty-first Chapter.}

\textbf{The Norms of the Class Invariants $f(\omega)$.}

\rule{0.3\linewidth}{0.4pt}
\end{center}

\begin{center}
§ 140. \gesperrt{Convergence of an Infinite Series.}
\end{center}

We consider an arbitrary quadratic form
\weq{1}{\psi(x,y) = Ax + 2Bxy + Cy^2} % [sic] source prints $Ax$; the quadratic form requires $Ax^2$
with negative discriminant
\weq{2}{-m = B^2 - AC,}
in which, for the present, the coefficients $A$, $B$, $C$ are not necessarily to be regarded as integers; they are only required to be real. If we take $x$, $y$ as rectangular coordinates in a plane, then to the integral values of $x$, $y$ there correspond the \gesperrt{lattice points}. The equation
\[\psi(x,y) = t\]
represents, for a constant $t$, an ellipse, and the number $Z(t)$ of those values of $\psi$, for \gesperrt{integral} $x$, $y$, which are smaller than $t$ is equal to the number of lattice points lying in the interior of this ellipse. According to Bd.\,II, § 194, (5), the limiting value of the ratio $Z(t):t$ for an infinitely large $t$ approaches the area of the ellipse $\psi = 1$, namely the value $\pi:\sqrt{m}$. Thus we have
\weq{3}{\mathrm{Lim}\, \frac{Z(t)}{t} = \frac{\pi}{\sqrt{m}}.}
If, then, in the infinite series
\weq{4}{S = \sum^{x,\,y} \frac{1}{\psi(x,y)^s}}
we let the numbers $x$, $y$ run through all integral values, with the exception of the combination $x=0$, $y=0$, then, according to Bd.\,II, § 196, 4, $S$ is convergent for $s>1$, and we have
\weq{5}{\mathop{\mathrm{Lim}}\limits_{s=1}\, (s-1)S = \frac{\pi}{\sqrt{m}}.}
